% ===========================================================================
% CHAPTER 9
% ===========================================================================

\chapter{Conclusion and Future Work}
\label{chap:conclusion-future-work}

\section{Conclusion}
\label{sec:conclusion}

SSL Guardian addresses the challenge of collecting, maintaining, analyzing, and
presenting large SSL/TLS certificate datasets through one integrated platform.
The developed system combines concurrent endpoint crawling, structured X.509
parsing, CT-assisted renewal and new-domain discovery, persistent certificate
storage, global and country-level analytics, application services, and an
interactive dashboard.

The acquisition framework uses persistent task states, atomic worker claims,
concurrent processing, retries, stale-task recovery, certificate parsing, scope
enrichment, and leaf classification. A separate IP-based crawler
extends the acquisition work to address-based measurement while remaining
isolated from the primary domain-based analytical dataset.

The CT pipeline supports continuous dataset development by observing
CT-derived domain sets, selecting known-domain renewal candidates, discovering
new domains, crawling their current endpoints, and incorporating successful
results. Analytical materialization follows completed update cycles, allowing
the dashboard to reflect newly acquired certificate information with minimal
operational delay while preserving a clear boundary from strict
event-by-event real-time processing.

The analytical platform covers validity, issuer and Certificate Authority
activity, signature and hash algorithms, key properties, SANs, public-key
reuse, geographic scope, trends, and security-oriented risk indicators.
Certificate Authority ranking and certificate-risk scoring provide transparent
comparative views, while the dashboard connects aggregate visualizations to
filtered lists and detailed certificate records. These scores and labels are
application-defined analytical aids and should not be interpreted as proof of
compromise or externally certified security ratings.

The project demonstrates a broad systems-engineering contribution to
certificate intelligence. Its principal strengths are the integration of
acquisition, CT-driven dataset maintenance, analytical processing, and
interactive exploration. The evaluation also identifies limitations in
production hardening, test automation, request isolation, benchmark coverage,
and analytical calibration. These boundaries define the present scope and
provide a clear direction for continued development.

\section{Future Work}
\label{sec:future-work}

Future engineering work should strengthen the platform through systematic
unit, integration, recovery, service-contract, and interface testing, supported
by reproducible datasets and explicit dependency versioning. Unified
configuration, externalized secrets, environment-specific settings,
authentication, role-based authorization, request-level scope isolation,
automated deployment, and operational monitoring would further support secure
and reproducible use beyond a locally operated research environment.

Larger deployments would benefit from distributed crawling, durable task
coordination, resource-aware scheduling, backpressure, and idempotent
processing. Incremental or distributed analytics could reduce repeated
full-dataset computation as the certificate dataset grows. Controlled
benchmarks should accompany these changes to evaluate acquisition throughput,
update latency, query responsiveness, and resource consumption across larger
datasets and analytical scopes.

Certificate intelligence may be extended through certificate-chain and trust
path analysis, deeper examination of certificate policies and name constraints,
and active OCSP or CRL status assessment. Retaining historical certificate
versions would enable lifecycle, renewal, key-rotation, issuer-migration, and
stale-deployment studies. Certificate Authority evaluation could consequently
incorporate transparent longitudinal measures, while ownership-aware analysis
could distinguish legitimate shared infrastructure from concerning public-key
reuse~\cite{junaid2025shared}.

Certificate Transparency processing could progress from scheduled new-domain
and renewal updates toward continuous, long-term ecosystem observation.
Durable per-log progress, event histories, and incremental correlation with
acquired certificates could improve renewal detection and reveal unusual
changes in issuers, public keys, SAN coverage, validity profiles, or issuance
frequency. Such anomalies should remain explainable investigative indicators.

Machine learning and artificial intelligence offer a complementary research
direction when sufficient history and reliable labels become available.
Certificate attributes, CT observations, issuer behaviour, renewal history,
cryptographic properties, quality findings, and analytical indicators could
support anomaly detection, risk estimation, suspicious-certificate
prioritization, and analyst triage. Models would require calibration,
interpretability, bias assessment, and concept-drift monitoring; they should
not be claimed to detect hacked websites directly or predict attacks.

Further extensions may integrate governed threat intelligence, including
trust-store, passive-DNS, reputation, and known malicious infrastructure data,
with clear provenance and licensing controls. IPv6
measurement, broader deployment environments, and richer longitudinal and
comparative visualizations would expand research utility. Collectively, these
directions reinforce the platform as a foundation for large-scale SSL/TLS
certificate intelligence, cybersecurity research, and advanced analytical
systems.
